\subsection{Мотивация регуляризации}

Всегда ли решение задачи регрессии единственно?
Вообще говоря, нет.
Так, если в выборке два признака будут линейно зависимы (и следовательно, ранг матрицы будет меньше $D$), то
гарантировано найдётся такой вектор весов $v$ что $(v,x_i) = 0~\forall x_i$.
В этом случае, если какой-то $w$ является решением оптимизационной задачи, то и $w+\alpha \nu$ тоже является
решением для любого $\alpha$.
То есть решение не только не обязано быть уникальным, так ещё может быть сколь угодно большим по модулю.
Это создаёт вычислительные трудности.
Малые погрешности признаков сильно возрастают при предсказании ответа, а в градиентном спуске накапливается
погрешность из-за операций со слишком большими числами.

Конечно, в жизни редко бывает так, что признаки строго линейно зависимы, а вот быть приближённо линейно зависимыми
они вполне могут быть.
Такая ситуация называется мультиколлинеарностью.
В этом случае у нас, всё равно, возникают проблемы, близкие к описанным выше.
Дело в том, что $X \nu \sim 0$ для вектора $v$, состоящего из коэффициентов приближённой линейной зависимости,
и, поэтому $X^T X \nu$, то есть матрица $X^T X$ снова будет близка к вырожденной.
Как и любая симметричная матрица, она диагонализируется в некотором ортонормированном базисе, и некоторые из
собственных значений $\lambda_i$ близки к нулю.
Если вектор $X^T y$ в выражении $(X^T X)^{-1} X^T y$ будет близким к соответствующему собственному вектору, то он
будет умножаться на $\frac{1}{\lambda_i}$, что опять же приведёт к появлению у $w$ очень больших по модулю
компонент (при этом $w$ ещё и будет вычислен с большой погрешностью из-за деления на маленькое число).
И, конечно же, все ошибки и весь шум из матрицы $X$ при вычислении $y \sim X w$ будут умножаться на эти большие и
неточные числа и возрастать во много-много раз, что приведёт к проблемам, от которых нас не спасёт никакое
сингулярное разложение.

Отметим, что в случае, когда несколько признаков линейно зависимы, веса $w_i$ при них теряют физический смысл.
Может оказаться, что вес признака, с ростом которого таргет, казалось бы, должен увеличиваться, станет отрицательным.
Это делает модель не только неточной, но и принципиально не интерпретируемой.
Вообще, неадекватность знаков или величины весов -- хорошее указание на мультиколлинеарность.

Для того чтобы справиться с этой проблемой, задачу обычно регуляризуют, то есть добавляют к ней дополнительное
ограничение на вектор весов.
Это ограничение можно, как и исходный loss, задавать по-разному, но, как правило, ничего сложнее, чем $L_1$- и
$L_2$-нормы, не требуется.

Вместо исходной задачи теперь предлагается решить такую:
\[ \min_w L(f,X,y) = \min_w (\abs{X w - y }^2_2 + \lambda \abs{w}_k^k)\]

$\lambda$ -- это очередной параметр, а $\abs{w}_k^k$ -- это либо $\abs{w}_2^2 = w_1^2 + \dots + w_D^2$, либо
$\abs{w}_1^1 = \abs{w_1} + \dots + \abs{w_D}$

Добавка $\lambda \abs{w}_k^k$ называется регуляризационным членом или регуляризатором, а число $\lambda$ --
коэффициентом регуляризации.
Коэффициент $\lambda$ является гиперпараметром модели и достаточно сильно влияет на качество итогового решения.
Его подбирают по логарифмической шкале (скажем, от 1e-2 до 1e+2), используя для сравнения моделей с разными
значениями $\lambda$ дополнительную валидационную выборку.
При этом качество модели с подобранным коэффициентом регуляризации уже проверяют на тестовой выборке, чтобы
исключить переобучение.

Отдельно надо договориться о том, что вес $w_0$ соответствующий отступу от начала координат (то есть признаку из
всех единичек), мы регуляризовать не будем, потому что это не имеет смысла: если даже все значения $y$ равномерно
велики, это не должно портить качество обучения.
Обычно это не отображают в формулах, но если придираться к деталям, то стоило бы написать сумму по всем весам,
кроме $w_0$: $\abs{w}_2^2 = \sum_{j=1}^D w_j^2$ и $\abs{w}_1 = \sum_{j=1}^D \abs{w_j}$

\subsection{Гребневая регрессия Ridge ($L_2$)}

В случае $L_2$-регуляризации решение задачи изменяется не очень сильно.

\subsubsection{Формула решения}
Продифференцировав новый loss по $w$, легко получить, что «точное» решение имеет вид:
\[ w = (X^T X + \lambda I)^{-1} X^T y\]

Отметим, что за этой формулой стоит и понятная численная интуиция: раз матрица $X^T X$ близка к вырожденной, то
обращать её сродни самоубийству.
Мы лучше слегка исказим её добавкой $\lambda I$, которая увеличит все собственные значения на $\lambda I$, отодвинув
их от нуля.
Да, аналитическое решение перестаёт быть «точным», но за счёт снижения численных проблем мы получим более
качественное решение, чем при использовании «точной» формулы.

В свою очередь, градиент функции потерь
\[ \min_w L(f_w,X,y) = (\abs{X w - y }^2 + \lambda \abs{w}^2)\]

по весам теперь выглядит так:
\[ \nabla_w L(f_w,X,y) = 2 X^T (X w - y) + 2 \lambda w \]

Подставив этот градиент в алгоритм стохастического градиентного спуска, мы получаем обновлённую версию приближенного
алгоритма, отличающуюся от старой только наличием дополнительного слагаемого.

\subsection{Lasso-регуляризация ($L_1$)}

Отметим, что $L_1$- и $L_2$-регуляризацию можно определять для любой функции потерь $L(w,X,y)$ (и не только в
задаче регрессии, а и, например, в задаче классификации тоже).
Новая функция потерь будет соответственно равна $L^\wave (w,X,y) = L(w,X,y) + \lambda \abs{w}_1$ или
$L^\wave (w,X,y) = L(w,X,y) + \lambda \abs{w}_2^2$

$L_2$-регуляризация работает прекрасно и используется в большинстве случаев, но есть одна полезная особенность
$L_1$-регуляризации: её применение приводит к тому, что у признаков, которые не оказывают большого влияния на
ответ, вес в результате оптимизации получается равным 0.
Это позволяет удобным образом удалять признаки, слабо влияющие на таргет.
Кроме того, это даёт возможность автоматически избавляться от признаков, участвующих в соотношениях приближённой
линейной зависимости, соответственно, спасает от проблем, связанных с мультиколлинеарностью, о которых мы писали выше.

\subsection{Вероятностная интерпретация}

\underline{Байесовский подход к регуляризации}.
Регуляризация имеет естественную интерпретацию в рамках байесовской статистики как априорное распределение на параметры.

\underline{Ridge-регрессия как апостериорное среднее}

\begin{theorem}[Байесовская интерпретация Ridge]
    Если:
    \begin{enumerate}
        \item Правдоподобие: $p(\bm{y} \mid \bm{X}, \bm{\beta}, \sigma^2) = \mathcal{N}(\bm{X}\bm{\beta}, \sigma^2 \bm{I}_n)$
        \item Априорное распределение: $p(\bm{\beta} \mid \tau^2) = \mathcal{N}(\bm{0}, \tau^2 \bm{I}_p)$
    \end{enumerate}
    Тогда апостериорное распределение:
    \[ p(\bm{\beta} \mid \bm{y}, \bm{X}, \sigma^2, \tau^2) = \mathcal{N}(\hat{\bm{\beta}}_{\text{Ridge}},
    \bm{\Sigma}_{\text{post}}) \]
    где:
    \begin{gather*}
        \hat{\bm{\beta}}_{\text{Ridge}} = \left(\frac{1}{\sigma^2}\bm{X}^\top\bm{X} +
                                             \frac{1}{\tau^2}\bm{I}_p\right)^{-1} \frac{1}{\sigma^2}\bm{X}^\top\bm{y}\\
        \bm{\Sigma}_{\text{post}} = \left(\frac{1}{\sigma^2}\bm{X}^\top\bm{X} + \frac{1}{\tau^2}\bm{I}_p\right)^{-1}\\
    \end{gather*}
\end{theorem}

\begin{proof}
    По теореме Байеса:
    \[ p(\bm{\beta} \mid \bm{y}, \bm{X}, \sigma^2, \tau^2) \propto p(\bm{y} \mid \bm{X}, \bm{\beta}, \sigma^2) \cdot
    p(\bm{\beta} \mid \tau^2)\]
    Логарифм апостериорного распределения:
    \begin{align*}
        \log p(\bm{\beta} \mid \cdot) &\propto -\frac{1}{2\sigma^2}\|\bm{y} - \bm{X}\bm{\beta}\|^2 -
        \frac{1}{2\tau^2}\|\bm{\beta}\|^2 \\
        &\propto -\frac{1}{2}\left[\frac{1}{\sigma^2}\|\bm{y} - \bm{X}\bm{\beta}\|^2 + \frac{1}{\tau^2}\|\bm{\beta}\|^2\right]
    \end{align*}
    Максимизация логарифма апостериорного распределения (MAP-оценка) эквивалентна минимизации:
    \[ \frac{1}{\sigma^2}\|\bm{y} - \bm{X}\bm{\beta}\|^2 + \frac{1}{\tau^2}\|\bm{\beta}\|^2 \],
    что совпадает с задачей Ridge-регрессии при $\lambda = \frac{\sigma^2}{\tau^2}$.
\end{proof}

\underline{Lasso как MAP-оценка}

\begin{theorem}[Байесовская интерпретация Lasso]
    Если:
    \begin{enumerate}
        \item Правдоподобие: $p(\bm{y} \mid \bm{X}, \bm{\beta}, \sigma^2) = \mathcal{N}(\bm{X}\bm{\beta}, \sigma^2 \bm{I}_n)$
        \item Априорное распределение: $p(\bm{\beta} \mid \lambda) = \prod_{j=1}^p \frac{\lambda}{2}
        e^{-\lambda|\beta_j|}$ (независимые распределения Лапласа)
    \end{enumerate}
    Тогда MAP-оценка совпадает с Lasso-решением:
    \[ \hat{\bm{\beta}}_{\text{MAP}} = \arg\max_{\bm{\beta}} p(\bm{\beta} \mid \bm{y}, \bm{X}) =
    \arg\min_{\bm{\beta}} \left\{ \|\bm{y} - \bm{X}\bm{\beta}\|^2 + 2\sigma^2\lambda \|\bm{\beta}\|_1 \right\} \]
\end{theorem}

\begin{proof}
    Априорное распределение Лапласа:
    \[ p(\bm{\beta} \mid \lambda) = \left(\frac{\lambda}{2}\right)^p e^{-\lambda \sum_{j=1}^p |\beta_j|} =
    \left(\frac{\lambda}{2}\right)^p e^{-\lambda \|\bm{\beta}\|_1} \]
    Логарифм апостериорного распределения:
    \begin{align*}
        \log p(\bm{\beta} \mid \bm{y}, \bm{X}) &\propto \log p(\bm{y} \mid \bm{X}, \bm{\beta}) + \log p(\bm{\beta} \mid \lambda) \\
        &\propto -\frac{1}{2\sigma^2}\|\bm{y} - \bm{X}\bm{\beta}\|^2 - \lambda \|\bm{\beta}\|_1
    \end{align*}
    Максимизация эквивалентна минимизации:
    \[ \frac{1}{2\sigma^2}\|\bm{y} - \bm{X}\bm{\beta}\|^2 + \lambda \|\bm{\beta}\|_1 \],
    что соответствует Lasso с $\alpha = \lambda\sigma^2$.
\end{proof}

\underline{Выводы из вероятностной интерпретации}

\begin{enumerate}
    \item Регуляризация соответствует введению априорных знаний о параметрах
    \item Гиперпараметр регуляризации $\lambda$ связан с параметрами априорного распределения
    \item Использование $L_2$-регуляризации предполагает, что все коэффициенты малы и имеют нулевое среднее.
    \item Использование $L_1$-регуляризации предполагает, что многие коэффициенты равны нулю (разреженность)
    \item Байесовский подход позволяет естественным образом оценивать неопределенность предсказаний
\end{enumerate}

\subsection{Elastic Net}

Можно пойти и далее, и соединить в себе оба регуляризационных члена:
\[ L_{EN} = \abs{Y - Xw}_2^2 + \lambda_1^2 \abs{w}_1 + \lambda_2^2 \abs{w}_2^2 \]
Что лучше работает?
Лучше работает кросс-валидация, но если выбирать из этих вариантов, то обычно $L_2$.

\subsection{Свойства регуляризаций}

\begin{table}[h]
    \centering
    \begin{tabular}{|l|l|l|}
        \hline
        \textbf{Метод} & \textbf{Априорное распределение}            & \textbf{Свойства}     \\ \hline
        МНК            & Неинформативное (равномерное)               & Нет регуляризации     \\ \hline
        Ridge          & Гауссово $\mathcal{N}(0, \tau^2 I)$      & Неразреженное решение \\ \hline
        Lasso          & Лапласовское $\text{Laplace}(0, 1/\lambda)$ & Разреженное решение   \\ \hline
        Elastic Net    & Комбинация Гаусса и Лапласа                 & Компромисс            \\ \hline
    \end{tabular}
    \caption{Байесовская интерпретация различных методов регуляризации}
    \label{tab:regul}
\end{table}